\subsection{Hak Asasi Manusia dalam UUD NRI Tahun 1945}

\subsubsection{Pasal 28F}
\begin{itemize}
\item Hak setiap orang untuk berkomunikasi dan memperoleh informasi guna mengembangkan pribadi dan lingkungan sosialnya.
\item Hak untuk mencari, memperoleh, memiliki, menyimpan, mengolah, dan menyampaikan informasi dengan menggunakan segala jenis saluran yang tersedia.
\end{itemize}

\subsubsection{Pasal 28G}
\begin{itemize}
\item Ayat (1): Hak atas perlindungan diri pribadi, keluarga, kehormatan, martabat, dan harta benda di bawah kekuasaannya, serta hak atas rasa aman dan perlindungan dari ancaman ketakutan untuk berbuat atau tidak berbuat sesuatu.
\item Ayat (2): Hak untuk bebas dari penyiksaan atau perlakuan yang merendahkan derajat martabat manusia dan hak memperoleh suaka politik dari negara lain.
\end{itemize}

\subsubsection{Pasal 28H}
\begin{itemize}
\item Ayat (1): Hak hidup sejahtera lahir dan batin, bertempat tinggal, mendapatkan lingkungan hidup yang baik dan sehat, serta berhak memperoleh pelayanan kesehatan.
\item Ayat (2): Hak mendapat kemudahan dan perlakuan khusus untuk memperoleh kesempatan dan manfaat yang sama guna mencapai persamaan dan keadilan.
\item Ayat (3): Hak atas jaminan sosial yang memungkinkan pengembangan dirinya secara utuh sebagai manusia yang bermartabat.
\item Ayat (4): Hak mempunyai hak milik pribadi dan hak milik tersebut tidak boleh diambil alih secara sewenang-wenang oleh siapa pun.
\end{itemize}

\subsubsection{Pasal 28I}
\begin{itemize}
\item Ayat (1): Hak asasi manusia yang tidak dapat dikurangi dalam keadaan apa pun (non-derogable rights), meliputi hak untuk hidup, hak tidak disiksa, hak kemerdekaan pikiran dan hati nurani, hak beragama, hak untuk tidak diperbudak, hak untuk diakui sebagai pribadi di hadapan hukum, dan hak untuk tidak dituntut atas dasar hukum yang berlaku surut.
\item Ayat (2): Hak bebas dari perlakuan yang bersifat diskriminatif atas dasar apa pun dan hak mendapatkan perlindungan terhadap perlakuan diskriminatif.
\item Ayat (3): Penghormatan terhadap identitas budaya dan hak masyarakat tradisional selaras dengan perkembangan zaman dan peradaban.
\item Ayat (4): Perlindungan, pemajuan, penegakan, dan pemenuhan hak asasi manusia adalah tanggung jawab negara, terutama pemerintah.
\end{itemize}

\subsubsection{Pasal 28J}
\begin{itemize}
\item Ayat (1): Setiap orang wajib tunduk kepada pembatasan yang ditetapkan dengan undang-undang dengan maksud semata-mata untuk menjamin pengakuan serta penghormatan atas hak dan kebebasan orang lain dan untuk memenuhi tuntutan yang adil sesuai pertimbangan moral, nilai-nilai agama, keamanan, dan ketertiban umum.
\end{itemize}

\subsection{Bela Negara}
\begin{itemize}
\item Pasal 27 Ayat (3): Setiap warga negara berhak dan wajib ikut serta dalam upaya pembelaan negara.
\item Pasal 30 Ayat (1): Tiap-tiap warga negara berhak dan wajib ikut serta dalam usaha pertahanan dan keamanan negara.
\end{itemize}

\subsection{Pendidikan dan Kebudayaan}

\subsubsection{Pasal 31}
\begin{itemize}
\item Ayat (1): Setiap warga negara berhak mendapat pendidikan.
\item Ayat (2): Setiap warga negara wajib mengikuti pendidikan dasar dan pemerintah wajib membiayainya.
\item Ayat (3): Pemerintah mengusahakan dan menyelenggarakan satu sistem pendidikan nasional yang meningkatkan keimanan dan ketakwaan serta akhlak mulia dalam rangka mencerdaskan kehidupan bangsa.
\item Ayat (4): Negara memprioritaskan anggaran pendidikan sekurang-kurangnya 20\% dari APBN dan APBD untuk memenuhi kebutuhan penyelenggaraan pendidikan nasional.
\item Ayat (5): Pemerintah memajukan ilmu pengetahuan dan teknologi dengan menjunjung tinggi nilai-nilai agama dan persatuan bangsa untuk kemajuan peradaban serta kesejahteraan umat manusia.
\end{itemize}

\subsubsection{Pasal 32}
\begin{itemize}
\item Ayat (1): Negara memajukan kebudayaan nasional Indonesia di tengah peradaban dunia dengan menjamin kebebasan masyarakat dalam memelihara dan mengembangkan nilai-nilai budayanya.
\item Ayat (2): Negara menghormati dan memelihara bahasa daerah sebagai kekayaan budaya nasional.
\end{itemize}

\subsection{Perekonomian Nasional dan Kesejahteraan Sosial}

\subsubsection{Pasal 33}
\begin{itemize}
\item Ayat (1): Perekonomian disusun sebagai usaha bersama berdasar atas asas kekeluargaan.
\item Ayat (2): Cabang-cabang produksi yang penting bagi negara dan yang menguasai hajat hidup orang banyak dikuasai oleh negara.
\item Ayat (3): Bumi dan air dan kekayaan alam yang terkandung di dalamnya dikuasai oleh negara dan dipergunakan untuk sebesar-besar kemakmuran rakyat.
\item Ayat (4): Perekonomian nasional diselenggarakan berdasar atas demokrasi ekonomi dengan prinsip kebersamaan, efisiensi berkeadilan, berkelanjutan, berwawasan lingkungan, kemandirian, serta dengan menjaga keseimbangan kemajuan dan kesatuan ekonomi nasional.
\item Ayat (5): Ketentuan lebih lanjut mengenai pelaksanaan pasal ini diatur dalam undang-undang.
\end{itemize}

\subsubsection{Pasal 34}
\begin{itemize}
\item Ayat (1): Fakir miskin dan anak-anak yang terlantar dipelihara oleh negara.
\item Ayat (2): Negara mengembangkan sistem jaminan sosial bagi seluruh rakyat dan memberdayakan masyarakat yang lemah dan tidak mampu sesuai dengan martabat kemanusiaan.
\item Ayat (3): Negara bertanggung jawab atas penyediaan fasilitas pelayanan kesehatan dan fasilitas pelayanan umum yang layak.
\item Ayat (4): Ketentuan lebih lanjut mengenai pelaksanaan pasal ini diatur dalam undang-undang.
\end{itemize}

\subsubsection{Contoh Hak Warga Negara Indonesia dalam Kehidupan Sehari-hari}
\begin{itemize}
\item Memeluk dan menjalankan ibadah sesuai dengan agama dan kepercayaannya.
\item Menyuarakan pendapat secara lisan maupun tulisan.
\item Mendapatkan pendidikan dan pengajaran.
\item Menikah dan berkeluarga.
\item Mendapatkan penghidupan yang layak.
\item Berorganisasi dan berkumpul.
\item Memiliki suatu barang atau properti yang legal.
\item Memiliki kedudukan yang sama di dalam hukum dan mendapat perlindungan hukum.
\end{itemize}

\subsubsection{Contoh Kewajiban Warga Negara Indonesia dalam Kehidupan Sehari-hari}
\begin{itemize}
\item Membayar pajak dan retribusi yang ditetapkan pemerintah.
\item Taat dan patuh pada peraturan hukum yang berlaku.
\item Menghormati orang lain.
\item Mengikuti pendidikan dasar.
\item Ikut menjaga sarana dan fasilitas umum.
\item Membuang sampah pada tempatnya.
\item Patuh kepada pembatasan atas hak kebebasan.
\end{itemize}

\subsection{Pentingnya Kesadaran Berkonstitusi}

\subsubsection{Pengertian Kesadaran Berkonstitusi}
\begin{itemize}
\item Meluasnya materi muatan UUD NRI Tahun 1945 setelah amandemen menuntut pemahaman warga negara terhadap konstitusi guna mewujudkan konstitusi yang hidup (living constitution).
\item Kesadaran berkonstitusi diartikan sebagai kualitas pribadi seseorang yang memancarkan wawasan, sikap, dan perilaku bermuatan cita-cita serta komitmen luhur kebangsaan dan kebernegaraan Indonesia.
\item Kesadaran berkonstitusi mengandung makna ketaatan pada aturan hukum sebagai aturan main (rule of the game) dalam kehidupan berbangsa dan bernegara.
\end{itemize}

\subsubsection{Penerapan Kesadaran Berkonstitusi}
\begin{itemize}
\item Budaya sadar berkonstitusi tercipta tidak hanya melalui pemahaman norma dasar, tetapi memerlukan pengalaman nyata dalam kehidupan bermasyarakat, berbangsa, dan bernegara.
\item Unsur-unsur kesadaran moral meliputi:
  \begin{enumerate}
  \item Perasaan wajib atau keharusan untuk melakukan tindakan bermoral sesuai dengan konstitusi.
  \item Rasional, yaitu bersifat objektif, berlaku umum, serta terbuka bagi pembenaran atau penyangkalan.
  \item Kebebasan atas kesadaran moral untuk menaati peraturan perundang-undangan dan ketentuan konstitusi.
  \end{enumerate}
\item Pembinaan kesadaran berkonstitusi dapat dilaksanakan melalui kegiatan kemasyarakatan seperti gotong royong, siskamling, dan rapat RT/RW.
\item Tanpa adanya kesadaran berkonstitusi, UUD NRI Tahun 1945 hanya akan menjadi dokumen tertulis tanpa aktualisasi dalam praktik kehidupan.
\end{itemize}

\subsubsection{Tingkatan Kesadaran Berkonstitusi}
\begin{itemize}
\item Kesadaran anomous: Kepatuhan terhadap ketentuan konstitusi yang tidak jelas dasar dan alasan atau orientasinya.
\item Kesadaran heteronomous: Kepatuhan terhadap ketentuan konstitusi yang berlandaskan motivasi beraneka ragam atau berganti-ganti.
\item Kesadaran socionomous: Kepatuhan terhadap ketentuan konstitusi yang berorientasi pada kiprah umum atau khalayak ramai.
\item Kesadaran autonomous: Kepatuhan terhadap ketentuan konstitusi yang didasari oleh konsep kesadaran dalam diri sendiri (tingkatan tertinggi).
\item Warga negara yang memiliki kesadaran berkonstitusi ditandai dengan kemampuan literasi konstitusi (constitutional literacy).
\end{itemize}

\subsubsection{Penerapan Hak dan Kewajiban Warga Negara}
\begin{itemize}
\item Pasal 27 ayat (1) mengenai persamaan kedudukan di dalam hukum:
  \begin{itemize}
  \item Hak: Memperoleh perlindungan dan persamaan hukum, seperti didampingi pengacara dan perlakuan adil dalam proses hukum.
  \item Kewajiban: Menghindari tindakan melanggar hukum, memenuhi panggilan saksi, serta menghormati lembaga dan aparat penegak hukum.
  \end{itemize}
\item Pasal 27 ayat (2) mengenai hak atas pekerjaan dan penghidupan yang layak:
  \begin{itemize}
  \item Hak: Memperoleh pekerjaan dan penghidupan yang layak, seperti mendaftar pekerjaan dan menerima upah.
  \item Kewajiban: Melaksanakan pekerjaan dengan penuh tanggung jawab, disiplin, jujur, dan memiliki etos kerja yang baik.
  \end{itemize}
\item Pasal 28 mengenai hak politik (berserikat, berkumpul, dan berpendapat):
  \begin{itemize}
  \item Hak: Mendirikan atau bergabung dengan organisasi, menyampaikan aspirasi, dan berkumpul untuk diskusi.
  \item Kewajiban: Mematuhi aturan hukum dan peraturan organisasi, bertanggung jawab atas aspirasi yang disampaikan, serta menyelesaikan sengketa secara damai.
  \end{itemize}
\item Pasal 28A sampai 28J mengenai hak atas HAM:
  \begin{itemize}
  \item Hak: Membentuk keluarga melalui perkawinan sah, kebebasan meyakini kepercayaan, dan memperoleh perlindungan hak.
  \item Kewajiban: Menghargai hak asasi orang lain, bertoleransi terhadap keragaman, dan mematuhi hukum yang berlaku.
  \end{itemize}
\item Pasal 29 mengenai kebebasan beragama:
  \begin{itemize}
  \item Hak: Beragama dan beribadah sesuai agama dan kepercayaannya.
  \item Kewajiban: Menjaga kerukunan antarumat beragama dan tidak memaksakan agama atau kepercayaan kepada orang lain.
  \end{itemize}
\item Perwujudan sikap kesadaran berkonstitusi meliputi:
  \begin{itemize}
  \item Kesadaran dan kesediaan mempertahankan serta mengisi kemerdekaan Indonesia.
  \item Pengakuan bahwa kemerdekaan merupakan rahmat Allah Yang Maha Kuasa.
  \item Kepekaan terhadap kewajiban negara dalam melindungi segenap bangsa, memajukan kesejahteraan umum, mencerdaskan kehidupan bangsa, dan ikut melaksanakan ketertiban dunia.
  \end{itemize}
\end{itemize}

\subsection{Perilaku Demokratis dalam Era Keterbukaan Informasi}

\subsubsection{Pengertian Demokrasi Pancasila}
\begin{itemize}
\item Demokrasi merupakan sistem yang mensyaratkan penghargaan terhadap harkat dan martabat manusia tanpa diskriminasi.
\item Nilai-nilai Demokrasi Pancasila bersumber dari Pembukaan UUD NRI Tahun 1945 Alinea IV, dengan esensi utama berupa kedaulatan rakyat.
\item Landasan konstitusional demokrasi tercantum pada Pasal 28 dan Pasal 28E Ayat (3) UUD NRI Tahun 1945 tentang kebebasan berserikat, berkumpul, dan mengeluarkan pendapat.
\item Kedisiplinan terhadap hukum dan etika menjadi fondasi utama dalam membangun demokrasi dan menciptakan ketahanan masyarakat.
\end{itemize}

\subsubsection{Pemahaman Demokrasi secara Normatif dan Empiris}
\begin{itemize}
\item Pemahaman secara Normatif: Demokrasi berada dalam kerangka berpikir konseptual sebagai landasan menata sistem pemerintahan untuk mewujudkan cita-cita dan tujuan negara dengan menempatkan peran penting rakyat.
\item Pemahaman secara Empiris: Demokrasi diwujudkan dalam praktik politik praktis dan penyelenggaraan pemerintahan yang dilandasi oleh persamaan hak, kewajiban, dan perlakuan adil bagi setiap warga negara.
\end{itemize}

\subsubsection{Hakikat dan Etimologi Demokrasi}
\begin{itemize}
\item Etimologi: Berasal dari bahasa Yunani "demos" (rakyat) dan "kratos" (kekuasaan atau berkuasa), yang berarti kekuasaan atau kedaulatan rakyat.
\item KBBI: Sistem pemerintahan yang seluruh rakyatnya turut serta memerintah dengan perantaraan wakilnya.
\item Tiga pilar utama penegak demokrasi:
  \begin{enumerate}
  \item Pemerintahan dari rakyat (government of the people).
  \item Pemerintahan oleh rakyat (government by the people).
  \item Pemerintahan untuk rakyat (government for the people).
  \end{enumerate}
\end{itemize}

\subsubsection{Klasifikasi Bentuk Demokrasi}
\begin{itemize}
\item Demokrasi secara Langsung:
  \begin{itemize}
  \item Pengertian: Memberikan kekuasaan kepada warga negara secara langsung tanpa perantara dalam pengambilan keputusan pemerintahan.
  \item Contoh: Pemilihan umum Presiden dan pemilukada langsung oleh rakyat.
  \end{itemize}
\item Demokrasi secara Tidak Langsung:
  \begin{itemize}
  \item Pengertian: Pengambilan keputusan pemerintahan dilakukan oleh wakil rakyat yang duduk dalam lembaga perwakilan rakyat.
  \item Contoh: Wakil rakyat yang menentukan kebijakan publik.
  \end{itemize}
\end{itemize}
